\section{Introduction}\label{sec:intro}

A developer who sees a qubit end in the wrong state asks which lines of the program determine it. For classical
programs a backward slice answers this~\cite{weiser1984program}, and programmers use slices to read and debug
code~\cite{weiser1982programmers, gallagher1991using}. Quantum programs need such help more. Their intermediate states
cannot be inspected without disturbing them, and a gate can change a qubit it does not target, as in phase kickback,
where a controlled gate changes the phase of its control without entangling the two
qubits~\cite{nielsen2010quantum}. Quantum developers debug
by printing, drawing circuits, simulating and reading code, and report that this does not
scale~\cite{zappin2026challenges}. Bug detectors such as LintQ and QChecker find known patterns by static analysis,
which reads a program without running it, % static-definition
but many faults are wrong outputs that no pattern describes~\cite{paltenghi2024lintq, zhao2023qchecker}.

Slicing has been proposed for quantum programs~\cite{yu2025qpdg} and named as an aid to understanding
them~\cite{murillo2025roadmap}, but no proposed dependence rule has been checked against what a quantum program computes. For comprehension this matters
most, because a slice that hides code able to change the qubit sends its reader away from the fault, and such slices
look right. The classical slice of the textbook program of Fig.~\ref{fig:bv} reproduces the faulty program's state
and leaves out the faulty line. Over the \numNPairs{} circuit--qubit pairs of our study, reading a controlled gate as
an assignment that leaves its control unchanged gives a wrong slice for \numCpdgUnsoundPct{} of the pairs, and
\numFdCpdgZMatchPct{} of those reproduce the whole program's outcome histogram, so a developer who checks a slice by
its outcomes would accept them.

This paper asks what a quantum slice must keep, and which dependence rule produces it. A slice of a set of qubits must
preserve their joint state for every input, and an operation is a dependence of them if a change right after it can
change that state. One rule meets both, the \emph{write-all rule}, under which a gate writes every qubit it touches,
its controls included. On qubits its slice is the backward light cone, whose soundness compilers already rely
on~\cite{murali2019triq, qiskit2025}. What has been missing is a proof that it keeps every dependence, so that no
other proposed kind of dependence can add one, and a measurement of what each alternative loses. We give the proof and
measure the losses on \numNCircuits{} benchmark circuits with an exact ground truth, \numFamCircuits{} generated
circuits, \numRqThreeUnits{} seeded-fault units and \numBugsFourQUnits{} real-fault units from
Bugs4Q~\cite{zhao2023bugs4q}. The study was pre-registered, but its confirmatory tests concern pruning, and the
comparisons of the other rules, added after the first results, are exploratory.

The paper makes four contributions. The first is a definition of what a slice of a set of qubits must preserve and of
dependence, which separates sound slices from slices that keep every dependence. The second is a proof that the slice
of the quantum dependence graph (QDG), defined by the write-all rule and extended along classical bits, is sound for
every input, with the induction mechanized in Isabelle/HOL, and keeps every dependence, with evidence that none of Qiskit's
standard two-qubit gates can safely write fewer qubits. The third shows that classical data, classical control and quantum data
dependence suffice, since apparent quantum control dependence is quantum data dependence and entanglement is neither
necessary nor sufficient for dependence. The fourth measures the alternatives. The classical PDG fails almost
entirely through its reading of the control. The QPDG of Yu and Zhao~\cite{yu2025qpdg} keeps \numYzLcExcessRatio{}
times as much code that cannot matter as the QDG, whose slices are on average \numLcExcess{} percentage points larger
than the smallest slice that keeps every dependence, a gap that grows with width. Pruning drops dependences. Slices
contain most faults but improve neither mutation-based fault ranking nor a language model's localization of faults.

The paper builds no slicer for a particular language. Its definitions are stated on the circuit a program builds.
Section~\ref{sec:lines} shows how a slice reaches source lines, and computing slices at scale on the source developers
write is a separate question. Section~\ref{sec:related} reviews related work, Section~\ref{sec:example} works through
Fig.~\ref{fig:bv}, Section~\ref{sec:core} defines the semantics and the rules, Sections~\ref{sec:design}
and~\ref{sec:results} present the study, and the last three sections give the implications, the threats to validity
and the conclusion. The data, scripts, protocol and proofs are in an anonymous replication package.\TODO{anonymous package link}
